\chapter{Prospective evaluation protocols}
\label{app:experiments}
\chapterpromise{Predefined statistical tests for posterior accuracy, calibration, irregular
spacing, pooling, finite latent groups, numerical integration, prediction and computational cost.}

The complete executable design matrices, seed construction, primary outcomes, confidence
interval methods, benchmark tuning and output filenames reside in one current experiment
specification. The experiments E00--E10 are \emph{planned}. No table or figure in this appendix
presents a result from those experiments. Earlier executed examples in
Chapters~\ref{ch:synthetic} and~\ref{ch:realdata} provide context under their original models,
not substitute confirmatory outcomes.

\section{Common statistical design}
All simulation comparisons use the same simulated sequence for every method in each replicate.
The independent generated dataset is the unit of replication. Posterior boundary calibration is
summarized by proper scoring rules, reliability curves with bin counts and dataset-level
confidence intervals. Boundary matching uses a maximum-cardinality, minimum-displacement
one-to-one assignment within the declared tolerance. Distinguish the posterior over segment
count from a conditional-on-$k$ boundary distribution. For real-data comparisons, coordinate
axes, train/validation/test splits and the reference label type must be verified before scoring.
Model selection and tuning do not use the test observations.

\begin{longtable}{@{}p{0.10\textwidth}p{0.26\textwidth}p{0.54\textwidth}@{}}
\caption{Planned experiments; no outcomes are supplied in this table.}\label{tab:book-planned-design}\\
\toprule
\textbf{ID} & \textbf{Research question} & \textbf{Core planned comparison or verification}\\
\midrule
\endfirsthead
\toprule
\textbf{ID} & \textbf{Research question} & \textbf{Core planned comparison or verification}\\
\midrule
\endhead
E00 & Exactness of inference & Exhaustive enumeration and direct segment integration for short sequences under several proper priors.\\
E01 & Recovery and calibration & Replicated Gaussian sequences with varied jump sizes, segment counts and posterior uncertainty targets.\\
E02 & Family generality & Gaussian, Poisson, Binomial, negative-Binomial and Beta-valued block models under appropriate descriptors.\\
E03 & Irregular design & Compare index-uniform, length-cohesion, gap-hazard and combined priors under controlled gaps.\\
E04 & Shared boundaries & Evaluate pooled and unpooled segmentation as subject count and boundary heterogeneity vary.\\
E05 & Finite-template clustering & Assess label-invariant clustering accuracy, objective monotonicity, restarts and template collapse.\\
E06 & Nonconjugate integration & Compare independent adaptive quadrature with identifiable Laplace, variational and EP approximations.\\
E07 & Model misspecification & Test heavy tails, count dispersion, missingness and dependence under declared alternative models.\\
E08 & Posterior prediction & Compare family-correct out-of-sample scores under disjoint held-out regions and fit uncertainty.\\
E09 & Applied comparisons & Use common coordinates, independently tuned comparators and genuinely external annotations where available.\\
E10 & Runtime and resource cost & Report stage-separated timings and peak memory as $n$ and $k_{\max}$ vary.\\
\bottomrule
\end{longtable}

\section{Primary reporting conventions}
For every experiment, report the complete design grid, number of planned and executed cells,
random-seed schedule, data and configuration hashes, failures and exclusion criteria.
Do not pool performance scores across incompatible response scales or observational
likelihoods. Numerical approximations must report both local evidence error and the effect on
global segmentation posteriors. Model-derived boundary agreement and externally labeled
boundary accuracy are different estimands and belong in separate tables. Only completed and
independently checked outcomes may replace the planned result slots in the single-column paper.
